% Nouveau EGA V (cote 25): archive cote 25, site page 56 (PDF page 57).
% typescript. Draft compared: none.
% AI-transcribed by Claude Opus 5.5 (claude-opus-5-5, adaptive thinking, effort high), Message Batches API, batch ega25b (msgbatch_01A7n5YuXubUPNK1zsepi8RG), 2026-11-13; checked against both supplied scan parts; see CHECK_DECISIONS.tsv.
\scanpage{25/56}
\mspage{51}

n  au lieu de \var{\xout{secti\unsure{o|n}}}{}{typewriter strike-out}{ega25_25_056-v0001} hyperplans.

\uline{Proposition} 9.3. \var{\xout{et}}{}{typewriter strike-out}{ega25_25_056-v0002} Supposons \var{\inserted[typed]{X géom irréd et}}{X géom irréd et}{insertion}{ega25_25_056-v0003} dim f(X) $\hand{\geq}$ \unsure{2}, \var{\xout{et}}{}{typewriter strike-out}{ega25_25_056-v0004} soit x$\hand{\in}$X(k), \var{\inserted{Soit $n\geq 2$}}{Soit $n\geq 2$}{insertion}{ega25_25_056-v0005} \var{\xout{et}}{}{typewriter strike-out}{ega25_25_056-v0006}
\var{\xout{\unsure{soit}}}{}{typewriter strike-out}{ega25_25_056-v0007} \var{\inserted[typed]{a\hand{)}}}{a\hand{)}}{insertion}{ega25_25_056-v0008} \var{\deleted{Alors pour n$\hand{\geq}$2}}{}{deletion}{ega25_25_056-v0009}, \var{\over{c}{C}}{C}{writing over}{ega25_25_056-v0010}onsidér\var{\over{a}{o}}{o}{writing over}{ega25_25_056-v0011}n\var{\over{t}{s}}{s}{writing over}{ega25_25_056-v0012} la famille \var{\inserted{(linéaire)}}{(linéaire)}{insertion}{ega25_25_056-v0013} des sections hyperplanes
de X relativement à $f_n=u_nf$ \uline{qui passent par} x, son élément
générique \var{\xout{est}}{}{typewriter strike-out}{ega25_25_056-v0014} définit une $Y^{(n)}_{\hand{\unsure{\xi|\zeta}}}$ qui est \uline{géométriquement} \var{\xout{intègre}}{}{typewriter strike-out}{ega25_25_056-v0015}irréd.

b) \var{\xout{Si f(X) n'est pas contenu dans un plan défini sur $\overset{\hand{\scriptscriptstyle -}}{k}$, et n$\hand{\geq}$2,}}{}{typewriter strike-out}{ega25_25_056-v0016}
\var{\xout{ou si n$\hand{\geq}$3,}}{}{typewriter strike-out}{ega25_25_056-v0017} Soit L une sous-variété linéaire de P passant par f(x) et
ne contenant pas f(X). \var{\deleted{Alors pour n$\hand{\geq}$2}}{}{deletion}{ega25_25_056-v0018}, \var{\over{c}{C}}{C}{writing over}{ega25_25_056-v0019}onsidér\var{\over{a}{o}}{o}{writing over}{ega25_25_056-v0020}n\var{\over{t}{s}}{s}{writing over}{ega25_25_056-v0021} la famille \var{\inserted[typed]{(linéaire)}}{(linéaire)}{insertion}{ega25_25_056-v0022} des
sections hyperplanes de\var{\xout{s}}{}{typewriter strike-out}{ega25_25_056-v0023} X relativement à \var{\xout{\unsure{l'immer}}}{}{typewriter strike-out}{ega25_25_056-v0024} $f_n$ \var{\xout{qui pa}}{}{typewriter strike-out}{ega25_25_056-v0025} , définis
par des hyperplans de P(n) \var{\xout{contenant}}{}{typewriter strike-out}{ega25_25_056-v0026} \uline{``tangents à L en x''} (i.e. définis
par des \var{\inserted{n-}}{n-}{insertion}{ega25_25_056-v0027}formes sur P qui \var{\inserted[typed]{sur L}}{sur L}{insertion}{ega25_25_056-v0028} s'annulent à l'ordre 2 au moins en x), son
membre générique \var{\deleted{définit}}{}{deletion}{ega25_25_056-v0029} \var{\inserted{est}}{est}{insertion}{ega25_25_056-v0030} une $Y^{(n)}_{\hand{\unsure{\xi|\zeta}}}$ géométriquement irréductible.

c) Supposons X lisse en x, \var{\xout{et f non}}{}{typewriter strike-out}{ega25_25_056-v0031} et \var{\deleted{soit}}{}{deletion}{ega25_25_056-v0032} n$\hand{\geq}$3, \var{\inserted{ou f(X) pas contenu dans un plan défini sur $\bar{k}$.}}{ou f(X) pas contenu dans un plan défini sur $\bar{k}$.}{insertion}{ega25_25_056-v0033} \var{\xout{alors}}{}{typewriter strike-out}{ega25_25_056-v0034} \var{\over{c}{C}}{C}{writing over}{ega25_25_056-v0035}onsidérons
la famille des sections hyperplanes \var{\inserted{\var{\deleted{\unsure{tangentielles}}}{}{deletion}{ega25_25_056-v0037}}}{\var{\deleted{\unsure{tangentielles}}}{}{deletion}{ega25_25_056-v0037}}{insertion}{ega25_25_056-v0036} \var{\xout{génériques}}{}{typewriter strike-out}{ega25_25_056-v0038} $Y^{(n)}_{\hand{\unsure{\xi|\zeta}}}$ de X relativement
à $f_n$ qui sont \hand{``}\uline{tangentes à X en x}\hand{''}\unsure{\var{\over{.}{,}}{,}{writing over}{ega25_25_056-v0039}} Alors le membre générique de ladite
définit une $Y^{(n)}_{\hand{\unsure{\xi|\zeta}}}$ géométriquement irréductible.

La démonstration est essentiellement triviale en termes des
critères de la fin du n\textsuperscript{o} précédent. Prenant un modèle affine de P
contenant f(x), on est ramené dans a) à trouver \var{\xout{une}}{}{typewriter strike-out}{ega25_25_056-v0040} trois polynômes
en les coordonnées $\unsure{T|X}_1,..,\unsure{T|X}_r$ , de degré \var{\over{\illegible{1}}{$\leq$}}{$\leq$}{writing over}{ega25_25_056-v0041}2, soient P,Q,R, tels que
Q(t)/P(t) et R(t)/P(t) soient alg indép sur k dans K (où K est le
corps des fonctions de X, et t = $(t_1,..,t_r)$ est le système d'élé-%
ments de K défini par les $T_i$ ); dans b), on exige de plus que P,Q,R
s'annulent à l'ordre 2 au moins sur L, qu'on peut d'ailleurs supposer
défini par des équations $T_1,..T_s = 0$ ; enfin dans \unsure{c|b}) c'est le même
pb, mais L étant l'image de l'espace tangent à X en x, et on se permet
\var{\inserted{éventuellement}}{éventuellement}{insertion}{ega25_25_056-v0042} de prendre P,Q,R \var{\xout{seulement}}{}{typewriter strike-out}{ega25_25_056-v0043} de degré \hand{3} , i.e. un peu plus de marge.

\mspage{56}
